\begin{textsample}{POS Dim 2 – mt – Score 26.00 – mt/mt\_em\_76.txt}  \label{ex:f2_pos_226}
Arquitetura curricular para EPT e possibilidades de distribuição de \textbf{carga} \textbf{horária} A proposta do Novo Ensino Médio apresenta Formação Geral Básica e possibilidades de flexibilização curricular, com o intuito de \textbf{atender} às várias \textbf{trajetórias} que são almejadas pelos estudantes.
A fim de possibilitar melhor compreensão de como poderá ser organizada a \textbf{oferta} da Educação Profissional Tecnológica no contexto do Novo Ensino Médio, apresenta -se abaixo a conceituação de algumas nomenclaturas relacionadas à EPT : 1.
Preparação básica para o trabalho : permite que o estudante desenvolva conhecimentos, atitudes, habilidades e comportamentos que garantam a preparação para o trabalho, traz orientação \textbf{profissional}, vocacional e leva o estudante a fazer escolhas das ocupações que melhor se adequem ao seu \textbf{perfil}.
A articulação aos eixos \textbf{estruturantes} pode ocorrer em módulo de preparação básica para o trabalho e/ou de forma transversal às unidades curriculares durante todo o desenvolvimento da trilha de aprofundamento.
Em anexo a este documento, apresentamos uma possibilidade na vertente transversal. 2.
Formação \textbf{Técnica} e Profissional : deve ser realizada na escola ou em instituição parceira que garanta ao estudante diploma \textbf{técnico} e \textbf{profissional}.
São atividades teóricas e/ou práticas de conhecimentos, habilidades e atitudes exigidas para o exercício das funções referentes a uma profissão.
As parcerias deverão ser firmadas pela \textbf{mantenedora} com instituições devidamente reconhecidas e legalizadas. 3.
Formação Inicial Continuada - FIC : são os \textbf{cursos} de Capacitação Profissional, de aperfeiçoamento e de atualização \textbf{profissional} de curta duração com \textbf{carga} \textbf{horária} que varia entre cento e \textbf{sessenta} horas ( 160 h ) e quatrocentos horas ( 400 h ).
Preparam os estudantes para a vida produtiva, social, bem como promovem a sua inserção e a reinserção no mundo do trabalho. 4.
Aprendizagem \textbf{profissional} : conforme \textbf{preconiza} na Resolução nº 03, de 21 de novembro 2018, art. 6º, inciso X, alínea C : > é a formação técnico-profissional compatível com o desenvolvimento físico, moral, psicológico e social do jovem de 14 a 24 anos de idade, previsto no § 4º, do art. 428, da Consolidação das Leis do Trabalho ( CLT ) e em \textbf{legislação} específica, caracterizada por atividades teóricas e práticas, metodicamente organizadas em tarefas de complexidade progressiva, conforme respectivo \textbf{perfil} \textbf{profissional} ( BRASIL, 2018 ).
Considerando a \textbf{oferta} do Ensino Médio em três mil horas ( 3.000 h ) e de acordo com a Lei nº 13.415/2017, apresentam -se abaixo quatro possibilidades de arquiteturas para a EPT.
Reforça -se que a definição sobre a(s) que melhor atende(m) às necessidades do Estado ocorrerá(ão) por meio de \textbf{regulamentação} da Seduc/MT e do CEE.
Essa proposta de arquitetura apresenta a seguinte flexibilização : BNCC, de mil e oitocentas horas ( 1.800 h ) ; Formação \textbf{Técnica} específica, a depender do \textbf{curso} \textbf{ofertado}, oitocentos ( 800 ), mil ( 1.000 ) ou mil e oitocentas horas ( 1.200 h ), em que se contempla a preparação básica para o trabalho organizada a partir dos eixos \textbf{estruturantes} propostos nos Referenciais Curriculares para a Elaboração de \textbf{Itinerários} \textbf{Formativos} ; Projeto de Vida, de cento e \textbf{sessenta} ( 160 h ) e \textbf{Eletivas}, de \textbf{duzentos} e quarenta ( 240 h ).
Se o \textbf{curso} \textbf{ofertado} for de oitocentas horas ( 800 h ), o Projeto de Vida e \textbf{Eletivas} terão a \textbf{carga} \textbf{horária} de acordo com o disposto no quadro apresentado acima.
Porém, se o \textbf{curso} \textbf{técnico} for de mil horas ( 1000 h ), tem -se cento e \textbf{sessenta} ( 160 ) de Projeto de Vida e quarenta horas ( 40 h ) para \textbf{Eletiva}.
E, ainda, se o \textbf{curso} \textbf{técnico} compreender mil e duzentas horas ( 1200 h ), o Projeto de Vida será trabalhado como componente curricular contemplado na matriz do \textbf{curso} \textbf{técnico}, considerando o interesse do estudante pela EPT.
Nessa proposta, a matrícula do aluno é única e as aulas são \textbf{ofertadas} em um único turno, sendo o \textbf{currículo} articulado sem sobreposição e organizado por meio de uma única trilha.
Essa proposta de arquitetura apresenta a mesma possibilidade de flexibilização do item 1, porém, cabe à \textbf{mantenedora} e à instituição parceira a verificação de viabilidade da articulação da proposta pedagógica do \textbf{curso} aos eixos \textbf{estruturantes} propostos nos Referenciais Curriculares para a Elaboração de \textbf{Itinerários} \textbf{Formativos}.
Para o aluno, são estabelecidas duas matrículas obrigatórias, uma para o Ensino Médio regular e outra para o \textbf{Curso} \textbf{Técnico} com \textbf{oferta} concomitante, que pode ser ministrado tanto na unidade escolar quanto na outra instituição.
Para essa forma de \textbf{oferta}, existe a possibilidade de construção de uma plataforma integrada para o monitoramento e o acompanhamento de todo o processo pedagógico pelas duas instituições ( mantenedora/instituição parceira ).
Essa proposta de arquitetura apresenta a seguinte flexibilização : BNCC, de mil e oitocentas horas ( 1.800 h ) ; preparação básica para o trabalho, organizada a partir dos eixos \textbf{estruturantes} propostos nos Referenciais Curriculares para a Elaboração de \textbf{Itinerários} \textbf{Formativos}, que perfaz um total de duzentas e quarenta horas ( 240 h ), mais a base teórica de quinhentas e \textbf{sessenta} horas ( 560 h ), envolvendo a prática da aprendizagem \textbf{profissional}, conforme \textbf{legislação} específica, podendo totalizar oitocentas ( 800 ), mil ( 1.000 ) ou mil e duzentas horas ( 1.200 h ) ; Projeto de Vida, de cento e \textbf{sessenta} horas ( 160 h ) ; \textbf{Eletivas}, de duzentas e quarenta horas ( 240 h ) ou quarenta horas ( 40 h ), a depender da \textbf{carga} \textbf{horária} envolvida na aprendizagem \textbf{profissional}.
A matrícula do estudante é única, parte do seu processo de aprendizagem acontece na prática do mundo do trabalho, com a possibilidade de receber o certificado de aprendiz ou \textbf{certificação} de \textbf{curso} \textbf{técnico}, mais certificado de aprendiz.
As \textbf{eletivas} podem ser \textbf{ofertadas} pelas unidades escolares, instituições parceiras ou, ainda, caracterizarem -se como um aprendizado em ação, no local de trabalho. 4 – Um conjunto de FICs articuladas entre si, organizadas no mesmo eixo, ou áreas tecnológicas ou por temas de uma área de potencial econômico Essa proposta de arquitetura apresenta a seguinte flexibilização : BNCC, de mil e oitocentas horas ( 1.800 h ) ; preparação básica para o trabalho, organizada a partir dos eixos \textbf{estruturantes} propostos nos Referenciais Curriculares para a Elaboração de \textbf{Itinerários} \textbf{Formativos}, que perfaz um total de duzentas e quarenta horas ( 240 h ) e \textbf{qualificação} por FICs, organizadas por eixos ou temas ; Projeto de vida, de cento e \textbf{sessenta} horas ( 160 h ) e \textbf{Eletivas}, com duzentas e quarenta horas ( 240 h ).
As FICs são de livre \textbf{oferta} e devem desenvolver competências de nível operacional e serem referenciadas na Classificação Brasileira de Ocupações ( CBO ) e/ou por orientações do sistema de ensino, com a possibilidade de \textbf{certificações} intermediárias.
A proposta pedagógica poderá se dar de forma integrada e/ou concomitante.
As FICs devem ser ministradas por \textbf{profissionais} habilitados, ou seja, bacharéis e tecnólogos ou por \textbf{profissionais} habilitados da base comum.
Estas devem ser articuladas entre si e ao mundo do trabalho, organizadas no mesmo eixo ou em áreas tecnológicas e/ou por temas de uma área de potencial econômico. \textbf{Certificação}, diploma e/ou histórico escolar Caracterizam -se como documentos legais que fazem parte do processo de avaliação.
Permitem o acompanhamento e evidenciam resultados das atividades pedagógicas que compõem o \textbf{currículo}.
Devem constar a descrição dos \textbf{Itinerários} \textbf{Formativos} ( Trilhas de Aprofundamento, \textbf{Eletivas} e Projeto de Vida ) vivenciados pelos estudantes, bem como as unidades curriculares e a \textbf{carga} \textbf{horária} \textbf{cursada}.
Com o intuito de melhor esclarecer as suas funcionalidades, apresenta -se abaixo a descrição e a função de cada um deles : - Diploma : deve evidenciar o título de \textbf{técnico} na \textbf{habilitação} \textbf{profissional} estudada, o eixo tecnológico e a área tecnológica de vinculação ; - Certificado : deve ser emitido somente no caso de o estudante ter \textbf{cursado} a Formação Geral Básica e o Itinerário \textbf{Formativo} completos ; - Histórico Escolar : evidencia o \textbf{perfil} \textbf{profissional} de conclusão do \textbf{curso}, as unidades curriculares \textbf{cursadas} e suas devidas \textbf{cargas} \textbf{horárias}, frequências e aproveitamento ; - \textbf{Certificação} Intermediária : quando a formação for estruturada e organizada em módulos, com terminalidades específicas e com o objetivo de qualificar o estudante para o trabalho.
Para além do exposto, o certificado, diploma e/ou histórico escolar devem descrever temas e projetos trabalhados, produtos realizados, habilidades desenvolvidas, além de participações em atividades relevantes, como representação estudantil, olimpíadas de conhecimento, campeonatos esportivos, espetáculos artísticos e culturais, congressos, gincanas, ações comunitárias, dentre outros.
Para efeito de registro, a responsabilidade pela emissão do certificado de conclusão do Ensino Médio fica a cargo da instituição de origem do estudante.

% matched lemmas: atender, carga, certificação, currículo, cursar, curso, duzentos, eletivo, estruturante, formativo, habilitação, horário, itinerário, legislação, mantenedor, oferta, ofertar, perfil, preconizar, profissional, qualificação, regulamentação, sessenta, trajetória, técnico
\end{textsample}
